\documentclass[10pt,a4paper]{amsart}
\usepackage[margin=19mm]{geometry}
\usepackage{amsmath,amssymb,mathtools}
\usepackage[scr=rsfs,cal=euler]{mathalfa}
\usepackage{xfrac}
\usepackage[unicode,hidelinks,bookmarks=false]{hyperref}
\newcommand{\ie}{i.e.\ }
\newcommand{\cf}{cf.\ }
\newcommand{\resp}{respectively\ }
\newcommand{\Subsection}{\subsection}
\providecommand{\nobreakdash}{\nobreak\hbox{-}}
\setlength{\emergencystretch}{2em}
\multlinegap0pt
\usepackage{bm}
\usepackage{bbm}
\usepackage{dsfont}
\usepackage{xspace}
\usepackage{cancel}
\usepackage{mathtools}
\usepackage{amsthm}
\makeatletter
\@ifundefined{theorem}{\newtheorem{theorem}{Theorem}}{}
\@ifundefined{lemma}{\newtheorem{lemma}[theorem]{Lemma}}{}
\makeatother
\def\text{\hbox}
\title[From length-preserving pushouts of graphs to one-surjective ]{From length-preserving pushouts of graphs to one-surjective pullbacks of graph algebras}
\author{Piotr M. Hajac, Mariusz Tobolski}
\date{Source: arXiv:2209.14470v6 (CC BY 4.0)}
\begin{document}
\maketitle

\noindent\emph{Source and licence.} From length-preserving pushouts of graphs to one-surjective pullbacks of graph algebras. Version arXiv:2209.14470v6 (CC BY 4.0), distributed under the Creative Commons Attribution 4.0 International licence, as recorded in the arXiv licence metadata for this exact version and shown on the abstract page https://arxiv.org/abs/2209.14470v6. Full rights notice: licenses/arxiv-2209.14470-CC-BY-4.0.txt.\par\smallskip

\noindent\textit{Editorial scope.} Counted unit: the Lemma whose source label is \texttt{answer} in the article (source0.tex lines 1235--1252, source label \texttt{answer}), delivered together with the article's own complete original proof of it (source0.tex lines 1253--1348, one \texttt{proof} environment of 96 lines). No mathematics is written, repaired or re-derived: statement and proof are the article's own text line for line. Required context retained uncounted: none. Every definition, display and label of those lines is retained unchanged. Omitted from this package and not redistributed: 1--1234, 1349--3462. Nothing inside a retained excerpt is omitted or altered: every retained body is delivered exactly as the article's lines are written, so no omission receipt of any kind accompanies this package. Citations: the retained excerpts carry no citation command at all, so no bibliography entry of the article is transcribed and no reference is redistributed. Reference adaptation: none was needed and none was made; every reference inside the delivered unit resolves inside it, so no unresolved reference or citation is produced. Printed slips kept literal: no printed word, symbol, hypothesis or number of the counted statement or of its proof was altered. Layout and class: the article's own class is \texttt{amsart} with packages including amssymb, amsmath, amsfonts, latexsym, color, amsthm, tikz, mathtools, times, calc, fontenc, inputenc, graphicx, xy; the wrapper is the lane's pinned \texttt{amsart} editorial wrapper at 10pt on A4, which restates the theorem-like environments the retained text needs and adds the article's own macro definitions that the retained text needs (\texttt{\textbackslash text}), with extra wrapper packages (bm, bbm, dsfont, xspace, cancel, mathtools). The delivered numbering is the unit's own. Assets: no retained span contains a figure environment, a graphics-inclusion command, a PDF-page inclusion, a table or a TikZ picture; no image file is copied into this package, the package declares no asset dependency and no image file is redistributed with it. Licence: version arXiv:2209.14470v6 of this article is distributed under the Creative Commons Attribution 4.0 International licence, as recorded in the arXiv licence metadata for this exact version and shown on the abstract page https://arxiv.org/abs/2209.14470v6. A full rights notice is delivered at licenses/arxiv-2209.14470-CC-BY-4.0.txt. Characters: every retained excerpt is pure ASCII, so the delivered engine, XeLaTeX with its default Latin Modern text font, reports no missing character.
\begin{lemma}\label{answer}
Let $E\stackrel{(f^0,f^1)}{\longleftarrow} G\stackrel{(g^0,g^1)}{\longrightarrow} F$ 
be graph homomorphisms.
Then there exists a natural map
\begin{equation*}
h\colon FP(E)\underset{FP(G)}{\amalg}FP(F)\longrightarrow FP\left(E\underset{G}{\amalg}F\right).
\end{equation*}
Moreover,  if the graph homomorphisms are such that
\vspace*{-2mm}
\begin{enumerate}
\item
both $f^0$ and $g^0$ are injective (vertex injectivity),
\item
$t_\amalg(x)=s_\amalg(y)\;\Rightarrow\;(x,y\in\iota_{E^1}(E^1)\text{ or }x,y\in\iota_{F^1}(F^1))$
(one color),
\end{enumerate}
then the natural map $h$ is bijective.
\end{lemma}

\begin{proof}
The natural map $FP(E)\amalg FP(F)\rightarrow FP(E\amalg_GF)$ exists because $FP$ is a functor. 
This map descends to $h$ by the universal
property of pushouts in the category of sets and maps. Since graph homomorphisms preserve the length
of paths, we obtain the decomposition
\begin{equation}
FP(E)\underset{FP(G)}{\amalg}FP(F)=\coprod_{n\in\mathbb{N}} FP_n(E)\underset{FP_n(G)}{\amalg}FP_n(F).
\end{equation}
Hence,  we can write $h$ on elements as follows:
\begin{equation}
h([p]):=\begin{cases}
[p] &\text{for } p\in \left(E^0\amalg F^0\right)\cup \left(E^1\amalg F^1\right),\\
([a_1],\ldots,[a_n]) &\text{for } p:=(a_1,\ldots,a_n)\in FP(E)\amalg FP(F),\\
&a_i\in E^1\amalg F^1,\;1\leq i\leq n,\;n\in\mathbb{N}\setminus\{0,1\}.
\end{cases}
\end{equation}

Now, using the two assumptions, we will define the inverse of~$h$. For starters, we put 
\begin{equation}\label{zeroone}
h^{-1}([p]):=[p]\quad
\text{when}\quad p\in \left(E^0\amalg F^0\right)\cup \left(E^1\amalg F^1\right).
\end{equation}
Next, let us take $([a_1],\ldots,[a_n])\in FP(E\amalg_GF)$
 with 
\begin{equation}
[a_i]\in E^1\underset{G^1}{\amalg}F^1,\; 1\leq i\leq n,\;
 n\in\mathbb{N}\setminus\{0,1\}. 
\end{equation}
Since $t_\amalg([a_i])=s_\amalg([a_{i+1}])$ for all $1\leq i\leq n-1$,
from the one-color condition we conclude that there exist $b_i\in E^1\amalg F^1$, $1\leq i\leq n$, such that $[b_i]=[a_i]$ for all $i$, and
\begin{equation}
\forall~1\leq i\leq n:~b_i\in E^1\quad\text{or}\quad\forall~1\leq i\leq n:~b_i\in F^1.
\end{equation}
Hence, $[t_E(b_i)]=[s_E(b_{i+1})]$ or $[t_F(b_i)]=[s_F(b_{i+1})]$. We can now apply the vertex-injectivity condition to infer that
$t_E(b_i)=s_E(b_{i+1})$ or $t_F(b_i)=s_F(b_{i+1})$. 

Furthermore, if $(b'_1,\ldots,b'_n)\in FP(E)$ or $(b'_1,\ldots,b'_n)\in FP(F)$, then 
\begin{equation}\label{implication}
\big{(}\forall~1\leq i\leq n:~[b_i]=[b'_i]\big{)}\quad \Longrightarrow\quad [(b_1,\ldots,b_n)]=[(b'_1,\ldots,b'_n)]
\end{equation}
Indeed, let $c^1_i,\ldots,c^{m}_i\in G^1$ be  sequences such that
\begin{gather}
\begin{matrix}
\forall\;1\leq i\leq n\colon\phantom{xxxxxxxxxx}\\
 f^1(c^1_i)=b_i\text{ and } g^1(c^1_i)=:b^1_i
\end{matrix}
\;,\ldots,\;
\begin{cases}
g^1(c^{m}_i)=:b^{m-1}_i\text{ and } f^1(c^{m}_i)=b'_i&\text{if $m$ is even},\\
f^1(c^{m}_i)=:b^{m-1}_i\text{ and } g^1(c^{m}_i)=b'_i&\text{if $m> 1$ is odd},
\end{cases}\nonumber\\
\text{or}\nonumber\\
\begin{matrix}
\forall\;1\leq i\leq n\colon\phantom{xxxxxxxxxx}\\
 g^1(c^1_i)=b_i\text{ and } f^1(c^1_i)=:b^1_i
\end{matrix}
\;,\ldots,\;
\begin{cases}
g^1(c^{m}_i)=:b^{m-1}_i\text{ and } f^1(c^{m}_i)=b'_i&\text{if $m>1$ is odd},\\
f^1(c^{m}_i)=:b^{m-1}_i\text{ and } g^1(c^{m}_i)=b'_i&\text{if $m$ is even}.
\end{cases}
\end{gather}
Note that, if the primed and unprimed edges are in the same graph, then the parity of all $c$-sequences is always even, and if they are in different graphs, then the parity is always odd. Therefore, although for each $i$ the length of the sequence $c^1_i,\ldots,c^{m}_i$ might be
different, we can always choose the longest such a sequence and extend shorter sequences by the constant extrapolation by an even number of elements. Next, we need to prove that $t_G(c^j_i)=s_G(c^j_{i+1})$. 
Depending on the parity of $j$ and the
above alternative between $E$ and $F$, we either have 
\begin{gather}
[f^1(c^j_i)]=[b_i]\quad\text{and}\quad [f^1(c^j_{i+1})]=[b_{i+1}]\nonumber\\
\text{or}\nonumber\\
[g^1(c^j_i)]=[b_i]\quad\text{and}\quad[g^1(c^j_{i+1})]=[b_{i+1}].
\end{gather}
 In the former case, by the vertex injectivity, we obtain
\begin{equation}\label{eq}
t_E(f^1(c^j_i))=t(b_i)\quad\text{and}\quad s_E(f^1(c^j_{i+1}))=s(b_{i+1}),
\end{equation}
where $t$ and $s$ are, respectively, $t_E$ and $s_E$, or $t_F$ and $s_F$, depending on whether
$b_i\in E^1$ or $b_i\in F^1$. It follows from \eqref{eq} that
\begin{equation}
f^0(t_G(c^j_i))=t_E(f^1(c^j_i))=t(b_i)=s(b_{i+1})=s_E(f^1(c^j_{i+1}))=f^0(s_G(c^j_{i+1})),
\end{equation}
so, from the injectivity of $f^0$, we get $t_G(c^j_i)=s_G(c^j_{i+1})$, as needed. In the latter case,
the reasoning is completely analogous but uses the injectivity of $g^0$ instead of the injectivity of~$f^0$.
Thus we have shown that $(c^j_1,\ldots,c^j_n)\in FP(G)$, $1\leq j\leq m$, is a sequence of paths
 implementing the desired equivalence relation between $(b_1,\ldots,b_n)$ and $(b'_1,\ldots,b'_n)$.

Finally, it follows from \eqref{implication} that we can define 
\begin{equation}
h^{-1}(([a_1],\ldots,[a_n])):=[(b_1,\ldots,b_n)]
\end{equation} 
Combining it with \eqref{zeroone} gives us
a map 
\begin{equation}
h^{-1}:FP\left(E\underset{G}{\amalg}F\right)\longrightarrow FP(E)\underset{FP(G)}{\amalg}FP(F),
\end{equation} 
which is, clearly, the inverse of~$h$. 
\end{proof}

\end{document}
